% =====================================================================
%  1bat-ud4-6.tex  —  SESSIÓ 6: Del paper al full de càlcul: primer contacte
%  (sense preàmbul: s'inclou des de main.tex)
% =====================================================================
\horatitol{Sessió 6 — Del paper al full de càlcul: primer contacte}%
{Reproduir amb l'ordinador el que ja sabem fer a mà: introduir dades i calcular suma, mitjana i recompte amb les primeres fórmules.}

\subsection{Idea clau}

Fins ara hem calculat els paràmetres \textbf{a mà} (sessions 2 a 5). Ara fem el salt
a l'eina que faràs servir com a tècnic: el \textbf{full de càlcul} (Excel o Google
Sheets). Avui hi dediquem un \textbf{primer contacte} ordenat, amb una idea clau: la
màquina ha de donar \textbf{el mateix} que el càlcul a mà; si coincideix, guanyem
confiança i confirmem que entenem què calcula.

\begin{idea}
\textbf{Primeres passes al full de càlcul}
\begin{itemize}[leftmargin=1.4em, itemsep=2pt]
  \item Una \textbf{cel·la} és cada caixa; es nomena per \textbf{columna+fila}: \texttt{A1},
        \texttt{B3}\dots Les dades en una columna ocupen, p.\,ex., de \texttt{A2} a
        \texttt{A11} (el \textbf{rang} \texttt{A2:A11}).
  \item Tota \textbf{fórmula} comença amb el signe \texttt{=}.
  \item \textbf{Suma:} \texttt{=SUMA(A2:A11)} \quad (en anglès \texttt{=SUM}).
  \item \textbf{Mitjana:} \texttt{=MITJANA(A2:A11)} \quad (en anglès \texttt{=AVERAGE}).
  \item \textbf{Recompte} (quantes dades hi ha): \texttt{=COMPTA(A2:A11)} \quad
        (en anglès \texttt{=COUNT}).
\end{itemize}
\textbf{Comprova sempre} que el resultat de la màquina coincideix amb el teu càlcul a
mà.
\end{idea}

\subsection{Exemples}

\begin{exemple}[introduir dades i sumar-les]
Volem portar al full de càlcul les hores de voluntariat $5,\ 3,\ 8,\ 5,\ 2,\ 5,\ 7,\ 4,\ 6$
(les de la sessió 2). Les escrivim a la columna A, de \texttt{A2} a \texttt{A10},
amb un títol a \texttt{A1}:

\begin{center}
\begin{tabular}{c|c}
 & \textbf{A} \\
\hline
1 & Hores \\
2 & 5 \\
3 & 3 \\
4 & 8 \\
\multicolumn{2}{c}{\quad$\vdots$\quad} \\
10 & 6 \\
\end{tabular}
\end{center}

A la cel·la \texttt{C1} escrivim \texttt{=SUMA(A2:A10)} i la màquina respon $45$.
És \emph{exactament} la suma que faríem a mà.
\end{exemple}

\begin{exemple}[la mitjana automàtica i la comprovació]
Amb les mateixes dades, a \texttt{C2} escrivim \texttt{=MITJANA(A2:A10)} i obtenim
$5$. A \texttt{C3} escrivim \texttt{=COMPTA(A2:A10)} i obtenim $9$ (hi ha $9$ dades).

\textbf{La comprovació clau:} a mà havíem calculat $\overline{x}=\dfrac{45}{9}=5$. El
full de càlcul dóna el mateix. A més, la mitjana es pot reconstruir com
\[
\frac{\texttt{=SUMA(A2:A10)}}{\texttt{=COMPTA(A2:A10)}}=\frac{45}{9}=5,
\]
cosa que confirma que la funció \texttt{MITJANA} fa, internament, «suma dividida pel
recompte».
\end{exemple}

\subsection{Exercicis resolts}

\begin{exresolt}[escriure les fórmules correctes]
Has introduït a la columna A, de \texttt{A2} a \texttt{A8}, els donatius diaris (€)
d'una entitat: $20$, $24$, $22$, $25$, $19$, $26$, $24$. Escriu les fórmules per
obtenir el total, la mitjana i el recompte, i digues quin resultat ha de donar
cadascuna.
\solucio
Les dades van de \texttt{A2} a \texttt{A8} (set valors). Les fórmules:
\begin{itemize}[leftmargin=1.4em, itemsep=1pt]
  \item Total: \texttt{=SUMA(A2:A8)} $\rightarrow 160$.
  \item Mitjana: \texttt{=MITJANA(A2:A8)} $\rightarrow \dfrac{160}{7}\approx 22,9$.
  \item Recompte: \texttt{=COMPTA(A2:A8)} $\rightarrow 7$.
\end{itemize}
Coincideix amb el càlcul a mà de la sessió 1 ($22,9$ kg/dia, allà eren kg): la
màquina confirma el resultat.
\resposta{\texttt{=SUMA(A2:A8)}$=160$; \texttt{=MITJANA(A2:A8)}$\approx 22,9$;
\texttt{=COMPTA(A2:A8)}$=7$.}
\end{exresolt}

\begin{exresolt}[detectar un rang mal seleccionat]
Un company volia la mitjana de $10$ dades (a \texttt{A2:A11}), però ha escrit
\texttt{=MITJANA(A1:A11)} i li surt un valor estrany. Què ha passat?
\solucio
Ha inclòs la cel·la \texttt{A1}, que conté el \textbf{títol} de la columna (un text,
no un número), o bé una cel·la de més. En seleccionar el rang \textbf{cal començar a
la primera dada} (\texttt{A2}), no a la capçalera. La fórmula correcta és
\texttt{=MITJANA(A2:A11)}. \textbf{Lliçó:} sempre cal comprovar que el rang agafa
\emph{totes} les dades i \emph{cap} cel·la de més (ni la capçalera).
\resposta{Ha inclòs la capçalera \texttt{A1} al rang; el correcte és
\texttt{=MITJANA(A2:A11)}.}
\end{exresolt}

\subsection{Exercicis per fer}

\noindent\textit{Practica directament a l'ordinador. Per a cada exercici, anota la
fórmula que has escrit i el resultat, i comprova'l amb el càlcul a mà.}

\begin{exercicis}
\begin{exlist}
  \item Introdueix a la columna A els infants atesos cada dia: $12$, $15$, $9$, $14$,
        $10$ (a \texttt{A2:A6}). Escriu les fórmules de \textbf{suma}, \textbf{mitjana}
        i \textbf{recompte} i anota'n els resultats.
  \item Comprova a mà la mitjana de l'exercici 1 (suma dividida pel recompte) i
        verifica que coincideix amb la del full de càlcul.
  \item Introdueix els visitants de $6$ tardes: $40$, $52$, $48$, $45$, $50$, $43$.
        Quina fórmula uses per saber \textbf{quantes tardes} hi ha? I per a la
        \textbf{mitjana} de visitants?
  \item \textbf{(Compte amb el rang)} Si les teves dades són a \texttt{A2:A6} però
        escrius \texttt{=SUMA(A2:A10)}, què passarà amb el resultat? I si escrius
        \texttt{=SUMA(A1:A6)} amb un títol a \texttt{A1}?
  \item \textbf{(Comprovació)} Introdueix les hores $4$, $6$, $8$, $6$, $6$ de la
        sessió 4. Calcula'n la mitjana amb el full de càlcul i comprova que dóna $6$,
        com vam obtenir a mà.
  \item \textbf{(Raonament)} Quin avantatge té calcular la mitjana amb el full de
        càlcul en comptes de fer-ho a mà? I quin \textbf{perill} hi ha si t'hi
        confies sense entendre què calcula? Explica-ho amb les teves paraules.
\end{exlist}
\end{exercicis}

% ---------------------------------------------------------------------
\begin{idea}
\textbf{Full de validació} — per a cada conjunt que portis al full de càlcul, omple
la fila i comprova que les dues columnes \emph{coincideixen}:
\begin{center}
\renewcommand{\arraystretch}{1.4}
\begin{tabular}{p{4.6cm} c c c}
\toprule
\textbf{Conjunt de dades} & \textbf{A mà} & \textbf{Full de càlcul} & \textbf{Coincideix?} \\
\midrule
Suma & \rule{1.6cm}{0.4pt} & \rule{1.6cm}{0.4pt} & \\
Mitjana & \rule{1.6cm}{0.4pt} & \rule{1.6cm}{0.4pt} & \\
Recompte & \rule{1.6cm}{0.4pt} & \rule{1.6cm}{0.4pt} & \\
\bottomrule
\end{tabular}
\end{center}
\end{idea}

\noindent\textbf{Atenció a la diversitat.} Si no tens experiència amb el full de
càlcul, segueix la pauta amb captures dels passos bàsics (escriure una dada,
seleccionar un rang, escriure una fórmula amb \texttt{=}); les parelles s'organitzen
perquè qui té més soltesa digital doni suport a qui en té menys.

% ---------------------------------------------------------------------
\fullsolucions{Sessió 6}
\begin{solucions}
\begin{sollist}
  \item \texttt{=SUMA(A2:A6)}$=60$; \texttt{=MITJANA(A2:A6)}$=12$;
        \texttt{=COMPTA(A2:A6)}$=5$.
  \item A mà: $\dfrac{60}{5}=12$. Coincideix amb el full de càlcul: correcte.
  \item Recompte: \texttt{=COMPTA(A2:A7)}$=6$ tardes. Mitjana:
        \texttt{=MITJANA(A2:A7)}$=\dfrac{278}{6}\approx 46,3$ visitants.
  \item Si el rang és massa gran (\texttt{A2:A10}) però només hi ha dades fins a
        \texttt{A6}, les cel·les buides no afecten la \emph{suma} (compten com a $0$),
        però sí poden afectar segons la funció; el més segur és ajustar el rang. Si
        inclous \texttt{A1} amb un títol (text), la \texttt{SUMA} l'ignora però una
        \texttt{MITJANA} o un \texttt{COMPTA} poden donar un resultat equivocat:
        \textbf{cal seleccionar exactament el rang de dades}.
  \item \texttt{=MITJANA(A2:A6)}$=6$. Coincideix amb el càlcul a mà de la sessió 4.
  \item Resposta oberta. Avantatge: estalvia molta feina i, si canvies una dada, tot
        es recalcula sol; és imprescindible amb llistes llargues. Perill: si t'hi
        confies sense entendre què calcula, no detectaràs un \textbf{rang mal
        seleccionat} ni un resultat absurd; per això cal saber comprovar-ho a mà.
\end{sollist}
\end{solucions}
